\documentclass[11pt]{article}
\usepackage{fontspec,amsmath,amssymb,amsthm,hyperref}
\usepackage[margin=1in]{geometry}
\setlength{\emergencystretch}{3em}
\title{Existence of the additive Langevin integral equation}
\author{LangevinLean mathematical exposition}
\date{}
\begin{document}\maketitle
\noindent Source module: \texttt{UnadjustedLangevin/LangevinIntegral.lean}. Compile this proof document with XeLaTeX or LuaLaTeX. The accompanying HTML pairs the mathematical argument with the exact Lean statements and records the scope of the formalization.


In a complete real normed space, a globally Lipschitz drift and a continuous additive path \(B\), \(B(0)=0\), determine a continuous solution of \(Y_t-x=-\int_0^tg(Y_s)ds+B_t\) on every finite nonnegative horizon. The construction is deterministic and later applies with the Banach space \(L^2(P)\).


\section{Solve on a short interval}


Remove the additive input by writing \(Y(t)=Z(t)+B(t)\). The new unknown satisfies \(Z(t)=x+\int_0^t f(s,Z(s))ds\), with \(f(s,z)=-g(z+B(s))\). This time-dependent drift is \(L\)-Lipschitz in \(z\) and continuous in time even though \(B\) need not be differentiable.



Bound \(\|B(t)\|\le C'\) on the compact interval and set \(G=\|g(x)\|+LC'\), radius \(a=2hG\), and drift bound \(M=G+La\). The condition \(Lh\le1/2\) proves \(Mh\le a\), giving the ball-invariance condition in \texttt{IsPicardLindelof}. The proof checks these bounds and invokes \texttt{ODE.FunSpace.exists\_isFixedPt\_next}.



The chosen fixed point gives a continuous \(Z\) with its integral identity. Add back \(B\), use \(B(0)=0\), and rewrite the negative drift integral to obtain \(Y(t)-x=-\int_0^tg(Y(s))ds+B(t)\). The solution is represented as a real-time function, with the required identity asserted on \([0,h]\).



\noindent Lean declarations: \texttt{exists\_langevin\_integral\_small}.

\section{splice langevin integral}


Use \(Y(t)\) up to time \(a\), and \(Z(t-a)\) afterward. The equality \(Z(0)=Y(a)\) makes the pasted path continuous at the join. Split \(\int_0^t=\int_0^a+\int_a^t\), translate the second integral, and add the two supplied equations; the terms \(B(a)\) cancel, giving the same equation on \([0,a+b]\).



\noindent Lean declarations: \texttt{splice\_langevin\_integral}.

\section{exists langevin integral}


Choose a positive segment length small enough that \(L\delta\le1/2\). Construct the first Picard segment, then repeatedly solve with initial state at the previous endpoint and shifted input \(B(a+s)-B(a)\). The splicing lemma combines finitely many segments covering the prescribed horizon. The local contraction restriction determines the construction mesh and imposes no restriction on the final horizon.



\noindent Lean declarations: \texttt{exists\_langevin\_integral}.

\section{exists lp langevin integral small}


Apply the Banach-space Picard theorem to \(H=L^2(P;K)\), \(g(X)=\operatorname{lpDrift}(\nabla U)(X)\). The lifted drift has the same Lipschitz constant as \(\nabla U\). Completeness of \(L^2\) and continuity of the supplied \(L^2\)-noise path are the exact fixed-point hypotheses.



\noindent Lean declarations: \texttt{exists\_lp\_langevin\_integral\_small}.

\section{exists lp langevin integral}


Use the finite-horizon Banach-space existence theorem with the lifted gradient drift. The conclusion is continuity in the \(L^2\) norm and an equality of \(L^2\) classes at every time in the horizon. The Bochner time integral therefore lives in the same Hilbert space; no pathwise differentiability of the random noise is asserted.



\noindent Lean declarations: \texttt{exists\_lp\_langevin\_integral}.

\end{document}
